\textbf{\textit{Investigué sobre los diferentes tipos de red:}}

\begin{itemize}
    \item \textbf{\textit{PAN (Personal Area Network):}} Es una red de corta distancia, usualmente de 10 m o menos, centrada en un individuo. Se utiliza para comunicar dispositivos personales como teléfonos, audífonos, etc. Incluye tecnologías como \textit{Bluetooth, Zigbee y NFC}.

    \item \textbf{\textit{MAN (Metropolitan Area Network):}} Es una red de alta velocidad que cubre una zona geográfica extensa. Suele utilizar microondas o fibra óptica para interconectar diversas LAN. Incluye tecnologías como \textit{redes de televisión por cable o anillos de fibra óptica municipales}.

    \item \textbf{\textit{CAN (Campus Area Network):}} Es una interconexión de varias redes de área local (LAN) dentro de un área geográfica como la UNAM. Es más grande que una LAN, pero más pequeña que una MAN. Se optimiza para el intercambio de recursos y la seguridad perimetral.

    \item \textbf{\textit{WLAN (Wireless Local Area Network):}} Es la variante inalámbrica de la LAN. Utiliza ondas de radio para la conectividad en áreas como oficinas o casas. Emplea el estándar IEEE 802.11, mejor conocido como Wi-Fi, y permite la movilidad dentro de un radio limitado por el punto de acceso (AP).

    \item \textbf{\textit{SAN (Storage Area Network):}} Es una red dedicada de alta velocidad y baja latencia que conecta servidores con dispositivos de almacenamiento. Opera a nivel de bloque, por lo que no sobrecarga la red LAN principal con tráfico.
\end{itemize}

\cite{tanenbaum2021networks}
